% GENERATED FILE. Edit canonical lessons, then run npm run generate:books.
% canonical-source-hash: fnv1a64:c89dd2f5afd0d8f4
% canonical-lessons: SA-C34-arm, SA-C34-neck, SA-C34-bone, SA-C34-blood, SA-C34-body

\chapter{More of the Body}
\label{ch:sa-limbs}
\hlchaptermodality{34}

\begin{chapteropening}
\textbf{By the end of this chapter:} \emph{I can name an arm, a neck, a bone, blood and the body itself in Sanskrit, and name the Latin word for bone that is a cousin of one of them.}

Complete SA-C34-body and demonstrate the chapter promise.
\end{chapteropening}

\section[bāhuḥ]{\sk{बाहुः} (bāhuḥ) — an arm}

\label{lesson:SA-C34-arm}



\begin{quote}
Before the new word: what did \emph{samudraḥ} mean?
\end{quote}

\subsection*{You'll want to know: \sk{बाहुः}}
\textbf{\sk{बाहुः}} (\emph{bāhuḥ}) — \textquotedblleft{}an arm\textquotedblright{}.

You already have \textbf{\sk{हस्तः}}, the hand. \textbf{\sk{बाहुः}} is the arm above it.

Its English cousin is \emph{bough} — the arm of a tree, which is exactly the picture Sanskrit had. Greek \emph{pēkhus}, \textquotedblleft{}forearm\textquotedblright{}, is the same word again, and it named the cubit that measured out temples and ships. Ends in \emph{-uḥ}, like \textbf{\sk{तरुः}} and \textbf{\sk{सेतुः}}.

\begin{grammarlens}[title={What you've built}]
A sixth part of the body, and the branch of a tree hiding inside it.
\end{grammarlens}

\subsection*{Your turn}
Say these aloud:

\begin{itemize}
\raggedright
  \item \emph{bāhuḥ}
  \item it once more, slowly
  \item \emph{samudraḥ}, then \emph{bāhuḥ}
\end{itemize}

\begin{tcolorbox}[breakable,colback=teal!4,colframe=teal!35!black,title={Before you move on}]
What does \sk{बाहुः} mean? (\textquotedblleft{}an arm\textquotedblright{}.) Say it once more.
\end{tcolorbox}
\section[grīvā]{\sk{ग्रीवा} (grīvā) — the neck}

\label{lesson:SA-C34-neck}



\begin{quote}
Before the new word: what did \emph{bāhuḥ} mean?
\end{quote}

\subsection*{You'll want to know: \sk{ग्रीवा}}
\textbf{\sk{ग्रीवा}} (\emph{grīvā}) — \textquotedblleft{}the neck\textquotedblright{}.

\textbf{\sk{ग्रीवा}} is the neck, and its family kept the swallowing that the old root began with.

Latin \emph{gula}, \textquotedblleft{}throat\textquotedblright{}, is a cousin, behind \emph{gullet}; Russian \emph{griva} is a mane, the hair along a neck; and Latin \emph{vorāre}, \textquotedblleft{}to devour\textquotedblright{}, belongs to the same wide family, giving \emph{devour}, \emph{voracious} and \emph{carnivore}. Ends in \textbf{-\sk{आ}}, like \textbf{\sk{तारा}} and \textbf{\sk{लता}}.

\begin{grammarlens}[title={What you've built}]
A neck, and a throat full of English words.
\end{grammarlens}

\subsection*{Your turn}
Say these aloud:

\begin{itemize}
\raggedright
  \item \emph{grīvā}
  \item it once more, slowly
  \item \emph{bāhuḥ}, then \emph{grīvā}
\end{itemize}

\begin{tcolorbox}[breakable,colback=teal!4,colframe=teal!35!black,title={Before you move on}]
What does \sk{ग्रीवा} mean? (\textquotedblleft{}the neck\textquotedblright{}.) Say it once more.
\end{tcolorbox}
\section[asthi]{\sk{अस्थि} (asthi) — a bone}

\label{lesson:SA-C34-bone}



\begin{quote}
Before the new word: what did \emph{grīvā} mean?
\end{quote}

\subsection*{You'll want to know: \sk{अस्थि}}
\textbf{\sk{अस्थि}} (\emph{asthi}) — \textquotedblleft{}a bone\textquotedblright{}.

\textbf{\sk{अस्थि}} is short, neuter and very old — the same shape as \emph{akṣi}, \textquotedblleft{}eye\textquotedblright{}, which you have already had.

Its cousins line up letter for letter: Latin \emph{os}, \emph{ossis}, behind \emph{osseous} and \emph{ossify}; Greek \emph{osteon}, behind \emph{osteopath}; Avestan \emph{ast}; Hittite \emph{hastai}. Hittite was written down three and a half thousand years ago, which makes this one of the oldest words you can still say out loud.

\begin{grammarlens}[title={What you've built}]
A bone, and one of the oldest words anyone has written down.
\end{grammarlens}

\subsection*{Your turn}
Say these aloud:

\begin{itemize}
\raggedright
  \item \emph{asthi}
  \item it once more, slowly
  \item \emph{grīvā}, then \emph{asthi}
\end{itemize}

\begin{tcolorbox}[breakable,colback=teal!4,colframe=teal!35!black,title={Before you move on}]
What does \sk{अस्थि} mean? (\textquotedblleft{}a bone\textquotedblright{}.) Say it once more.
\end{tcolorbox}
\section[raktam]{\sk{रक्तम्} (raktam) — blood}

\label{lesson:SA-C34-blood}



\begin{quote}
Before the new word: what did \emph{asthi} mean?
\end{quote}

\subsection*{You'll want to know: \sk{रक्तम्}}
\textbf{\sk{रक्तम्}} (\emph{raktam}) — \textquotedblleft{}blood\textquotedblright{}.

\textbf{\sk{रक्तम्}} began as a description rather than a name: it means \textquotedblleft{}reddened, dyed\textquotedblright{}, from a root \emph{rañj}, \textquotedblleft{}to colour\textquotedblright{}. Sanskrit names blood by its colour and lets the colour word do the work.

The same root gives \textbf{\sk{राग}} (\emph{rāga}), which is a colour, a passion and a mode of music all at once — a thing that has been dyed a particular way. Neuter \textbf{-\sk{अम्}}, like \textbf{\sk{हिमम्}}.

\begin{grammarlens}[title={What you've built}]
Blood, named for its colour rather than for what it is.
\end{grammarlens}

\subsection*{Your turn}
Say these aloud:

\begin{itemize}
\raggedright
  \item \emph{raktam}
  \item it once more, slowly
  \item \emph{asthi}, then \emph{raktam}
\end{itemize}

\begin{tcolorbox}[breakable,colback=teal!4,colframe=teal!35!black,title={Before you move on}]
What does \sk{रक्तम्} mean? (\textquotedblleft{}blood\textquotedblright{}.) Say it once more.
\end{tcolorbox}
\section[dehaḥ]{\sk{देहः} (dehaḥ) — the body}

\label{lesson:SA-C34-body}



\begin{quote}
Before the last word of the chapter: what did \emph{raktam} mean, and what did \emph{bāhuḥ} mean?
\end{quote}

\subsection*{You'll want to know: \sk{देहः}}
\textbf{\sk{देहः}} (\emph{dehaḥ}) — \textquotedblleft{}the body\textquotedblright{}.

\textbf{\sk{देहः}} is the whole of it — the thing the other four belong to. Its root means \textquotedblleft{}to knead, to shape, to plaster\textquotedblright{}, so a body is a moulded thing, formed the way clay is formed.

English \emph{dough} is the plain cousin. Latin \emph{fingere}, \textquotedblleft{}to shape\textquotedblright{}, gives \emph{figure}, \emph{fiction} and \emph{effigy}. And an Old Persian relative, a wall of kneaded clay, sits inside \emph{paradise}: \emph{pairi-daēza}, a walled garden.

\begin{grammarlens}[title={What you've built}]
Ten parts of the body between this chapter and the earlier one, and a single word for all of them.
\end{grammarlens}

\subsection*{Your turn}
Say these aloud:

\begin{itemize}
\raggedright
  \item \emph{dehaḥ}
  \item it once more, slowly
  \item all five in order, then say what \emph{paradise} was built out of
\end{itemize}

\begin{tcolorbox}[breakable,colback=teal!4,colframe=teal!35!black,title={Before you move on}]
What does \sk{देहः} mean? (\textquotedblleft{}the body\textquotedblright{}.) Say it once more.
\end{tcolorbox}
